\section{负多重覆叠问题}
\label{sec:negative-multiple-covers}
\stfhIndex{I210}{fu^duo^chong^fu^die}{负多重覆叠（negative multiple cover）}

本节先说明：若要\emph{通过扰动 $M$ 上的殆复结构}来应用横截性定理，那么在某处单射这一假设是必要的。先从复几何中下面两个著名定理开始。以
$\mathcal O(a)$ 表示复直线丛中次数为 $a$ 的全纯截面之层。

\begin{theorem}\label{thm:grothendieck-splitting-cp1}
$\mathbb{CP}^1$ 上的每个全纯向量丛 $E$ 都分裂为
\[
 E\cong\mathcal O(a_1)\oplus\cdots\oplus\mathcal O(a_n),
 \qquad n=\operatorname{rank}(E),
\]
而整数集 $\{a_1,\ldots,a_n\}$ 与分解的选取无关。
\end{theorem}

\begin{theorem}\label{thm:fredholm-regular-splitting-criterion}
设 $M$ 是复流形。若 $u:\mathbb{CP}^1\to M$ 且
\[
 u^*TM\cong\mathcal O(a_1)\oplus\cdots\oplus\mathcal O(a_n),
\]
则 $u$ 弗雷德霍姆正则，当且仅当对所有 $1\leq i\leq n$ 都有
$a_i\geq-1$。
\end{theorem}

\begin{example}\label{ex:negative-multiple-cover-blowup}
设 $M=\mathbb{CP}^2$，取 $p\in M$，并令 $\widetilde M$ 是 $M$ 在
$p$ 处的爆破：
\[
 \widetilde M\xrightarrow{\ \pi\ }M.
\]
则 $\pi^{-1}(p)=E$ 是一条例外球面。固定嵌入
$u:S^2\to\widetilde M$，使 $u(S^2)=E$。若 $NE$ 是 $E$ 在
$\widetilde M$ 中的法丛，则 $NE$ 同构于复直线丛
\[
 \mathcal L:=\{(l,v)\in\mathbb{CP}^1\times\mathbb C^2\mid v\in l\},
\]
即 $NE\cong\mathcal O(-1)$。由定理 10.6.2 中 $u$ 的正则性判据，
$u$ 是弗雷德霍姆正则的全纯球面。

另一方面，考虑映射
$\phi:\mathbb{CP}^1\to\mathbb{CP}^1$，在 $\mathbb C$ 上由
$\phi(z)=z^2$ 给出，并考虑复合 $u\circ\phi$。于是
\[
 (u\circ\phi)^*T\widetilde M
 \cong T\mathbb{CP}^1\oplus(u\circ\phi)^*NE
 \cong\mathcal O(2)\oplus\mathcal O(-2).
\]
更一般地，$E$ 的重数为 $a$ 的多重覆叠有同构于
$\mathcal O(-a)$ 的法丛，其中 $a\geq2$。所以再由定理 10.6.2，这个多重覆叠并非弗雷德霍姆正则。

现在断言，这个多重覆叠不能被 $\widetilde M$ 上标准复结构的任意小扰动消去。这是因为，参数化例外球面
$E$ 的单覆叠
$u:S^2\to\widetilde M$ 对标准复结构是横截的，因而在把复结构小扰动为殆复结构后必然仍然存在。它的指标为
\[
 c_1(E)+2=2-1+2=3.
\]
因此，重数为 $a$ 的多重覆叠映射之指标为
$2-a+2=4-a$。例如，当 $a>4$ 时，指标变为负数；\emph{倘若模空间是横截的}，相应模空间便必须为空。这就完成了断言的证明。

负多重覆叠问题促使我们寻找研究横截性的新办法，允许采用比扰动殆复结构 $J$ 更一般的扰动。后者只依赖\emph{靶}流形
$M$，因而无法破坏\emph{定义域}所诱导的任何对称性。
\end{example}
